\section*{Capítulo \ref{ch:b1:atmospheric-circulation} --- A circulação geral da atmosfera}
% values: calc/ch28.py

\begin{solution}{exo:b1:atmospheric-circulation:1}
(a)~$f = 2\Omega\sin\varphi$: $0$ no equador, \qty{7.29e-5}{s^{-1}} a \ang{30},
\qty{1.031e-4}{s^{-1}} a \ang{45} e \qty{1.458e-4}{s^{-1}} no polo. (b)~$fV = \num{1.031e-4} \times 10 = \qty{1.03e-3}{m.s^{-2}}$, cerca de \num{1.05e-4} de $g$. (c)~No Hemisfério Norte o desvio é para a
direita do movimento: um vento que sopra para o norte é empurrado para leste.
\end{solution}

\begin{solution}{exo:b1:atmospheric-circulation:2}
(a)~$|\nabla p|/\rho = (400/\num{4e5})/1.225 = \qty{8.16e-4}{m.s^{-2}}$.
(b)~$10/\num{8.16e-4} = \qty{12250}{s}$, \qty{3.4}{h}. (c)~A aceleração de Coriolis de um vento de
\qty{10}{m/s} a \ang{45} N é \qty{1.03e-3}{m.s^{-2}}: as duas são do mesmo tamanho (razão $0.79$).
Em poucas horas depois de começar a se mover, o ar sente uma força de Coriolis capaz de equilibrar o
gradiente de pressão: o vento se estabiliza perto da velocidade geostrófica, aqui $\num{8.16e-4}/\num{1.031e-4} = \qty{7.9}{m/s}$,
em vez de acelerar descendo o gradiente.
\end{solution}

\begin{solution}{exo:b1:atmospheric-circulation:3}
(a)~$f = 2 \times \num{7.292e-5} \times \sin\ang{40} = \qty{9.374e-5}{s^{-1}}$;
$|\vec v_g| = (400/\num{2.5e5})/(1.225 \times \num{9.374e-5}) = \qty{13.9}{m/s}$, ou
$13.93/0.5144 = \qty{27.1}{kn}$. (b)~No Hemisfério Norte o vento deixa a baixa à sua esquerda: com a
baixa ao sul, sopra para oeste, vindo do leste (\omterm{def:b1:spherical-astronomy:horizontal}{azimute} \ang{90}). (c)~A haste aponta para leste a
partir do círculo da estação, para dentro do vento; \qty{27.1}{kn} arredonda para \qty{25}{kn}, duas
penas inteiras e meia pena, desenhadas no lado sul da haste, em direção à baixa.
\end{solution}

\begin{solution}{exo:b1:atmospheric-circulation:4}
(a)~\ang{15} N: \omterm{def:b1:atmospheric-circulation:winds}{alísios} de nordeste; \ang{15} S: \omterm{def:b1:atmospheric-circulation:winds}{alísios} de sudeste; \ang{45} N e \ang{45} S: \omterm{def:b1:atmospheric-circulation:winds}{ventos de
oeste}. (b)~A \ang{15} o ar de superfície da \omterm{def:b1:atmospheric-circulation:cells}{célula de Hadley} volta em direção ao equador e é desviado
para oeste pela força de Coriolis (para a direita no norte, para a esquerda no sul): de leste, $u$
negativo. A \ang{45} o ar de superfície da \omterm{def:b1:atmospheric-circulation:cells}{célula de Ferrel} se move para o polo e é desviado para leste:
de oeste, $u$ positivo. (c)~Das três \omterm{def:b1:spherical-astronomy:geographic}{latitudes}, a pressão é mais alta a \ang{30} N, \qty{1016.8}{hPa}: a
\omterm{def:b1:atmospheric-circulation:cells}{alta subtropical} fica perto de \ang{30} N (o perfil completo põe seu máximo perto de \ang{35} N).
\end{solution}

\begin{solution}{exo:b1:atmospheric-circulation:5}
(a)~$|\nabla p|/\rho = (200/\num{5e5})/1.225 = \qty{3.27e-4}{m.s^{-2}}$, dividido por $f$:
\qty{2.6}{m/s} a \ang{60}, \qty{2.9}{m/s} a \ang{50}, \qty{4.5}{m/s} a \ang{30}, \qty{6.5}{m/s} a
\ang{20}, \qty{12.9}{m/s} a \ang{10} e \qty{25.7}{m/s} a \ang{5}. (b)~A força de Coriolis gira um vento
em um tempo da ordem de $1/f$: \qty{2.2}{h} a \ang{60}, \qty{11}{h} a \ang{10}, \qty{22}{h} a \ang{5}, e
infinito no equador. Perto do equador o campo de pressão muda, e o atrito age, muito antes de o vento
poder ser girado ao longo das isóbaras: o equilíbrio nunca se estabelece, e o ar flui em parte descendo
o gradiente, a velocidades modestas.
\end{solution}

\begin{solution}{exo:b1:atmospheric-circulation:6}
(a)~$u = 464.6\,\sin^2\varphi/\cos\varphi$: \qty{14.2}{m/s} a \ang{10}, \qty{32.2}{m/s} a \ang{15},
\qty{57.8}{m/s} a \ang{20}, \qty{134.1}{m/s} a \ang{30}. (b)~$u = 464.6\,(\cos^2\ang{10} - \cos^2\varphi)/\cos\varphi$:
\qty{17.7}{m/s} a \ang{15}, \qty{42.9}{m/s} a \ang{20}, \qty{117.9}{m/s} a \ang{30}. (c)~O ar que
começa mais longe do equador começa mais perto do eixo, com menos momento angular sobrando: chega com
um vento de oeste mais fraco em cada \omterm{def:b1:spherical-astronomy:geographic}{latitude}. Espalhar o ramo ascendente baixa o jato previsto, mas
nem de longe o suficiente para atingir os \qty{40}{m/s} observados.
\end{solution}

\begin{solution}{exo:b1:atmospheric-circulation:7}
(a)~O vento sopra para nordeste. De costas para ele, de frente para o nordeste, a baixa fica à
esquerda: a noroeste. (b)~O vento sopra para sudeste; no Hemisfério Sul a baixa fica à direita de quem
está de costas para o vento: a sudoeste. (c)~O vento de superfície é desviado das isóbaras em direção à
baixa, isto é, no sentido anti-horário em relação ao \omterm{def:b1:atmospheric-circulation:geostrophic}{vento geostrófico} no Hemisfério Norte. O \omterm{def:b1:atmospheric-circulation:geostrophic}{vento
geostrófico}, e a baixa em ângulo reto à sua esquerda, são portanto girados no sentido horário em
relação à resposta de (a): a baixa fica entre noroeste e norte, no \omterm{def:b1:spherical-astronomy:horizontal}{azimute} $\ang{315} + \alpha$ para um
ângulo de cruzamento $\alpha$.
\end{solution}

\begin{solution}{exo:b1:atmospheric-circulation:8}
(a)~A área entre o equador e \ang{30} N é $2\pi a_\oplus^2\sin\ang{30} =
\qty{1.275e14}{m^2}$; a massa acima de um metro quadrado entre os dois níveis é
$\Delta p/g = \num{9e4}/9.81 = \qty{9174}{kg}$; o total é
\qty{1.170e18}{kg}. (b)~$\num{1.170e18}/\num{8.7e10} = \qty{1.34e7}{s}$, cerca de
\num{156} dias, cinco meses. (c)~$\num{1.170e18}/\num{21.4e10} = \qty{5.47e6}{s}$,
cerca de \num{63} dias, dois meses.
\end{solution}

\begin{solution}{exo:b1:atmospheric-circulation:9}
(a)~$\num{0.01} \times \num{86400} = \qty{864}{m}$ por dia. (b)~$0.864 \times 9.8 =
\qty{8.5}{K}$ por dia. (c)~O ar esfria emitindo radiação infravermelha: a atmosfera perde mais energia
por radiação do que absorve, e a descida do ar, aquecendo-o por compressão, compensa a perda.
(d)~O ar descendente aquece e seca em relação à saturação, suas nuvens se evaporam e a chuva não pode
se formar: as \omterm{def:b1:atmospheric-circulation:cells}{altas subtropicais} são cinturões de céu claro e pouca chuva, onde se formam desertos
quentes sobre a terra.
\end{solution}

\begin{solution}{exo:b1:atmospheric-circulation:10}
(a)~$V^2 + frV - fr|\vec v_g| = 0$ com $fr = \qty{55.85}{m/s}$:
$V = \tfrac12\bigl(-55.85 + \sqrt{55.85^2 + 4 \times 55.85 \times 15}\bigr) =
\qty{12.3}{m/s}$, \qty{82}{\percent} do \omterm{def:b1:atmospheric-circulation:geostrophic}{vento geostrófico}. (b)~$V^2 - frV +
fr|\vec v_g| = 0$ tem raízes reais apenas se $(fr)^2 \geq 4fr|\vec v_g|$, isto é
$|\vec v_g| \leq fr/4 = \qty{14.0}{m/s}$. (c)~Perto do centro de uma alta, onde $r$ é
pequeno, o limite $fr/4$ é pequeno: o gradiente de pressão deve ser fraco, as isóbaras bem
afastadas e os ventos leves. Uma depressão não tem tal limite: suas isóbaras podem se apertar
perto do centro, com ventos fortes.
\end{solution}

\begin{solution}{exo:b1:atmospheric-circulation:11}
(a)~Para $\varphi$ pequeno, $\sin\varphi \approx \varphi$ e $\cos\varphi \approx 1$,
logo $u \approx \Omega a_\oplus\varphi^2$. (b)~$\ang{10} = \qty{0.1745}{rad}$:
$464.6 \times 0.1745^2 = \qty{14.15}{m/s}$, contra \qty{14.23}{m/s} exatos,
\qty{0.5}{\percent} a menos. (c)~Resolvendo $464.6\sin^2\varphi/\cos\varphi = 27.0$
numericamente dá $\varphi = \ang{13.7}$. (d)~O anel já coincide com o jato anual médio real a menos de
\ang{14}: a conservação do momento angular descreve o ramo superior da \omterm{def:b1:atmospheric-circulation:cells}{célula de Hadley} apenas nos
trópicos profundos; a partir de cerca de \ang{15} em direção ao polo, os redemoinhos assumem.
\end{solution}

\begin{solution}{exo:b1:atmospheric-circulation:12}
(a)~Média $31.7/12 = \ang{2.6}$ N; amplitude de \ang{12.8} S a \ang{10.9} N, isto é \ang{23.7}.
(b)~O máximo ao norte vem em setembro, três meses depois do solstício de junho; o máximo ao sul em
fevereiro, dois meses depois do solstício de dezembro. O atraso vem do armazenamento de calor, sobretudo
na \omterm{def:b1:fluid-envelopes-structure:ocean-layers}{camada de mistura} do oceano (\cref{ch:b1:fluid-envelopes-structure}), que aquece e esfria devagar.
(c)~A média zonal reúne todas as \omterm{def:b1:spherical-astronomy:geographic}{longitudes}: se sobre os continentes a faixa de chuva oscila bastante
para o hemisfério de verão, a média de todas as \omterm{def:b1:spherical-astronomy:geographic}{longitudes} pode cruzar o equador embora a faixa sobre
um oceano não o faça.
\end{solution}

\begin{solution}{pb:b1:atmospheric-circulation:1}
\begin{enumerate}
  \item $\rho = p/(R_dT) = \num{100600}/(287.05 \times 283.15) =
        \qty{1.238}{kg.m^{-3}}$.
  \item $|\nabla p| = \qty{400}{Pa}/\qty{2.5e5}{m} = \qty{1.6e-3}{Pa.m^{-1}}$.
  \item $|\nabla p|/\rho = \num{1.6e-3}/1.238 = \qty{1.293e-3}{m.s^{-2}}$.
  \item $\Omega = 2\pi/\num{86164.2} = \qty{7.292e-5}{rad.s^{-1}}$;
        $f = 2\Omega\sin\ang{50} = \qty{1.117e-4}{s^{-1}}$.
  \item $|\vec v_g| = \num{1.293e-3}/\num{1.117e-4} = \qty{11.57}{m/s}$.
  \item $11.57/0.514444 = \qty{22.5}{kn}$; $11.57 \times 3.6 = \qty{41.7}{km/h}$.
  \item O tempo é $|\vec v_g|/(|\nabla p|/\rho)$ e $|\vec v_g| = (|\nabla
        p|/\rho)/f$, logo ele é $1/f = \qty{8951}{s}$, \qty{2.49}{h}; e $2\pi/f =
        \qty{15.6}{h}$. Padrões que duram um dia ou mais são portanto próximos do equilíbrio.
  \item A \omterm{def:b1:atmospheric-circulation:pgf}{força do gradiente de pressão} aponta para o norte, em direção à baixa. O \omterm{def:b1:atmospheric-circulation:geostrophic}{vento geostrófico}
        sopra ao longo das isóbaras com a baixa à sua esquerda: para leste, vindo do oeste, \omterm{def:b1:spherical-astronomy:horizontal}{azimute}
        \ang{270}: um vento de oeste.
  \item De costas para um vento de oeste, de frente para leste, a esquerda é o norte, onde fica a
        baixa: a lei é satisfeita.
  \item A haste aponta para oeste a partir do círculo da estação, na direção de onde o vento vem;
        \qty{22.5}{kn} arredonda para \qty{20}{kn}, duas penas inteiras; elas ficam do lado norte da
        haste, em direção à baixa.
  \item A \ang{50} S, $f < 0$: o vento mantém a baixa à sua direita, de modo que sopra para oeste,
        vindo do leste.
  \item $f(\ang{30}) = \qty{7.292e-5}{s^{-1}}$: \qty{17.7}{m/s}; $f(\ang{60}) =
        \qty{1.263e-4}{s^{-1}}$: \qty{10.2}{m/s}.
  \item $f(\ang{5}) = \qty{1.271e-5}{s^{-1}}$ dá \qty{101.7}{m/s}, um vento nunca observado com
        gradientes desse tipo. O tempo de giro $1/f$ é de \qty{22}{h}: o equilíbrio não tem tempo de se
        estabelecer, e a fórmula não se aplica.
  \item $9.0/11.57 = 0.778$; ao longo das isóbaras $9.0\cos\ang{25} = \qty{8.16}{m/s}$, através delas
        em direção à baixa $9.0\sin\ang{25} = \qty{3.80}{m/s}$.
  \item O \omterm{def:b1:atmospheric-circulation:geostrophic}{vento geostrófico} sopra para o \omterm{def:b1:spherical-astronomy:horizontal}{azimute} \ang{90}; girado \ang{25} em direção à baixa, no norte,
        o vento de superfície sopra para \ang{65}, vindo de \ang{245}, o oés-sudoeste.
  \item $\tan\alpha = 1/(f\tau)$ dá $\tau = 1/(\num{1.117e-4}\tan\ang{25}) =
        \qty{1.92e4}{s}$, \qty{5.3}{h}; o modelo prevê $11.57\cos\ang{25} =
        \qty{10.5}{m/s}$, mais que os \qty{9.0}{m/s} observados: o atrito real não é exatamente linear no
        vento, e freia mais que o modelo nesse ângulo.
  \item Através de um cilindro de raio $R = \qty{500}{km}$ e altura $h =
        \qty{1}{km}$: $\dot m = 2\pi R h\rho u_r = 2\pi \times \num{5e5} \times 1000
        \times 1.238 \times 3.80 = \qty{1.48e10}{kg/s}$.
  \item $w = \dot m/(\rho\pi R^2) = 2hu_r/R = 2 \times 1000 \times 3.80/\num{5e5} =
        \qty{0.0152}{m/s}$, cerca de \qty{1310}{m} por dia.
  \item $1.31 \times 9.8 = \qty{12.9}{K}$ por dia: o ar logo atinge a saturação, e a depressão se enche
        de nuvens e chuva. Sob um \omterm{def:b1:atmospheric-circulation:cyclone}{anticiclone} o ar da superfície escoa para fora, o ar acima desce e
        aquece, e o céu se limpa.
  \item $\Omega a_\oplus = \num{7.292e-5} \times \num{6.371e6} = \qty{464.6}{m/s}$ no equador;
        $464.6\cos\ang{30} = \qty{402.3}{m/s}$ a \ang{30} N.
  \item $u = 464.6\,\sin^2\varphi/\cos\varphi$: \qty{57.8}{m/s} a \ang{20} N e \qty{134.1}{m/s} a
        \ang{30} N.
  \item $u = 464.6\,(\cos^2\ang{10} - \cos^2\ang{30})/\cos\ang{30} =
        \qty{117.9}{m/s}$.
  \item $134.1/42.3 = 3.17$. Os redemoinhos das \omterm{def:b1:spherical-astronomy:geographic}{latitudes} médias, depressões e \omterm{def:b1:atmospheric-circulation:cyclone}{anticiclones}, tiram
        momento angular do ramo superior e o levam para outro lugar; e o ar da célula sobe sobre um
        conjunto de \omterm{def:b1:spherical-astronomy:geographic}{latitudes}, em vez de ficar em repouso no equador. A média zonal também faz a média de
        um jato sinuoso sobre todas as \omterm{def:b1:spherical-astronomy:geographic}{longitudes}.
  \item $M/M_0 = (402.3 + 42.3)\cos\ang{30}/464.6 = 0.829$: o ar perdeu \qty{17}{\percent} do momento
        angular que manteria como um anel livre.
  \item A velocidade geostrófica em P é \qty{11.6}{m/s}, de oeste; o vento medido no solo é
        \qty{9.0}{m/s}, cerca de quatro quintos dela, e girado \ang{25} em direção à baixa pelo atrito.
\end{enumerate}
\end{solution}
